\documentclass[final]{agujournal2019}
\usepackage{url}
\usepackage{lineno}
\usepackage{soul}

\usepackage{xcolor}
\definecolor{default-linkcolor}{HTML}{A50000}
\definecolor{default-filecolor}{HTML}{A50000}
\definecolor{default-citecolor}{HTML}{4077C0}
\definecolor{default-urlcolor}{HTML}{4077C0}

\usepackage{hyperref}
\hypersetup{
  hidelinks,
  breaklinks=true
}


\providecommand{\tightlist}{\setlength{\itemsep}{0pt}\setlength{\parskip}{0pt}}
\providecommand{\listfigurename}{List of Figures}
\providecommand{\listtablename}{List of Tables}
\newcounter{none}
\setcounter{none}{0}

\makeatletter
\@ifundefined{pandocbounded}{
  \usepackage{graphicx}
  \newcommand{\pandocbounded}[1]{\setbox0=\hbox{#1}\ifdim\wd0>\linewidth\resizebox{\linewidth}{!}{#1}\else#1\fi}
}{}
\makeatother

\renewcommand{\thefigure}{S\arabic{figure}}
\renewcommand{\thetable}{S\arabic{table}}
\renewcommand{\thesection}{S\arabic{section}}
\renewcommand{\theequation}{S\arabic{equation}}

\usepackage{longtable,booktabs,array}
\usepackage{calc}
\usepackage{etoolbox}
\makeatletter
\patchcmd\longtable{\par}{\if@noskipsec\mbox{}\fi\par}{}{}
\makeatother
\IfFileExists{footnotehyper.sty}{\usepackage{footnotehyper}}{\usepackage{footnote}}
\makesavenoteenv{longtable}

\NewDocumentCommand\citeproctext{}{}
\NewDocumentCommand\citeproc{mm}{\begingroup\def\citeproctext{#2}\cite{#1}\endgroup}
\makeatletter
 \let\@cite@ofmt\@firstofone
 \def\@biblabel#1{}
 \def\@cite#1#2{{#1\if@tempswa , #2\fi}}
\makeatother

\newlength{\cslhangindent}
\setlength{\cslhangindent}{1.5em}
\newlength{\csllabelwidth}
\setlength{\csllabelwidth}{3em}
\newenvironment{CSLReferences}[2]
 {\begin{list}{}{
  \setlength{\itemindent}{0pt}
  \setlength{\leftmargin}{0pt}
  \setlength{\parsep}{0pt}
  \ifodd #1
   \setlength{\leftmargin}{\cslhangindent}
   \setlength{\itemindent}{-\cslhangindent}
  \fi
  \setlength{\itemsep}{#2\baselineskip}
  }
  \item\relax}
 {\end{list}}
\usepackage{calc}
\newcommand{\CSLBlock}[1]{#1\hfill\break}
\newcommand{\CSLLeftMargin}[1]{\parbox[t]{\csllabelwidth}{#1}}
\newcommand{\CSLRightInline}[1]{\parbox[t]{\linewidth - \csllabelwidth}{#1}\break}
\newcommand{\CSLIndent}[1]{\hspace{\cslhangindent}#1}

\usepackage{amsmath}
\usepackage{amssymb}
\usepackage{amsfonts}
\makeatletter
\@ifpackageloaded{subcaption}{}{\usepackage{subcaption}}
\@ifpackageloaded{caption}{}{\usepackage{caption}}
\captionsetup[subfigure]{margin=0.5em}
\AtBeginDocument{%
\renewcommand*\figurename{Figure}
\renewcommand*\tablename{Table}
}
\AtBeginDocument{%
\renewcommand*\listfigurename{List of Figures}
\renewcommand*\listtablename{List of Tables}
}
\newcounter{pandoccrossref@subfigures@footnote@counter}
\newenvironment{pandoccrossrefsubfigures}{%
\setcounter{pandoccrossref@subfigures@footnote@counter}{0}
\begin{figure}\centering%
\gdef\global@pandoccrossref@subfigures@footnotes{}%
\DeclareRobustCommand{\footnote}[1]{\footnotemark%
\stepcounter{pandoccrossref@subfigures@footnote@counter}%
\ifx\global@pandoccrossref@subfigures@footnotes\empty%
\gdef\global@pandoccrossref@subfigures@footnotes{{##1}}%
\else%
\g@addto@macro\global@pandoccrossref@subfigures@footnotes{, {##1}}%
\fi}}%
{\end{figure}%
\addtocounter{footnote}{-\value{pandoccrossref@subfigures@footnote@counter}}
\@for\f:=\global@pandoccrossref@subfigures@footnotes\do{\stepcounter{footnote}\footnotetext{\f}}%
\gdef\global@pandoccrossref@subfigures@footnotes{}}
\@ifpackageloaded{float}{}{\usepackage{float}}
\floatstyle{ruled}
\@ifundefined{c@chapter}{\newfloat{codelisting}{h}{lop}}{\newfloat{codelisting}{h}{lop}[chapter]}
\floatname{codelisting}{Listing}
\newcommand*\listoflistings{\listof{codelisting}{List of Listings}}
\makeatother

\journalname{JGR: Solid Earth}

\makeatletter
\def\titlepageheadlinefont{\fontsize{9}{11}\selectfont}

\def\ps@headings{
  \def\@oddfoot{\centerline{\small --\the\c@page--}}
  \let\@evenfoot\@oddfoot
  \def\@oddhead{}
  \let\@evenhead\@oddhead
}
\ps@headings

\def\ps@front{
  \def\@oddhead{
    \vbox to 0pt{
      \vskip-30pt
      \centerline{\titlepageheadlinefont JGR: Solid Earth}
      \vss
    }
  }
  \let\@evenhead\@oddhead
  \def\@oddfoot{\centerline{\small --\the\c@page--}}
  \let\@evenfoot\@oddfoot
}

\let\oldtitle\title
\renewcommand{\title}[1]{\oldtitle{#1}\thispagestyle{front}}
\makeatother

\begin{document}


\title{Supplementary Information for "Beyond Equilibrium II: Sluggish
Ringwoodite Kinetics Amplify 660 km Discontinuity Depressions and
Independently Promote Slab Stagnation"}

\authors{Buchanan Kerswell\affil{1}, John Wheeler\affil{1}, Rene
Gassmöller\affil{2}, J. Huw Davies\affil{3}, Isabel
Papanagnou\affil{4}, Sanne Cottaar\affil{4}, David P. Dobson\affil{5}}

\affiliation{1}{Department of Earth, Oceans and Ecological Sciences,
University of Liverpool, Liverpool L69 3BX, UK}
\affiliation{2}{GEOMAR Helmholtz Centre for Ocean Research Kiel, Kiel,
Germany}
\affiliation{3}{School of Earth and Environmental Sciences, Cardiff
University, Park Place, Cardiff, CF10 3AT, UK}
\affiliation{4}{Bullard Laboratories, Department of Earth Sciences,
University of Cambridge, Cambridge CB3 0EZ, UK}
\affiliation{5}{Department of Earth Sciences, University College London,
London WC1E 6BS, UK}

\correspondingauthor{Buchanan Kerswell}{b.kerswell@liverpool.ac.uk}

\section*{Contents of this file}\label{sec:contents}
\addcontentsline{toc}{section}{Contents of this file}

\begin{enumerate}
\def\labelenumi{\arabic{enumi}.}
\tightlist
\item
  Figures \ref{fig:material-property-profiles} to
  \ref{fig:slab-scen11-set2}
\item
  Table~\ref{tbl:depth-profile-summary}
\end{enumerate}

\clearpage

\section*{Introduction}\label{sec:introduction}
\addcontentsline{toc}{section}{Introduction}

This supplementary information provides 18 figures and one table that
support the material model, kinetic calibration, and quantitative
simulation results described in the main text.

Figures \ref{fig:material-property-profiles} and
\ref{fig:thermodynamic-property-profiles} show the reference material
and thermodynamic properties used to initialize and drive the
compressible mantle flow simulations. These profiles give the density,
thermal expansivity, compressibility, and specific heat capacity of each
mantle phase across the transition zone, along with the Gibbs free
energy, entropy, and volume changes that govern the 410, 520, and 660
phase transitions.

Figure~\ref{fig:regression} presents the Monte Carlo regression used to
calibrate the ringwoodite decomposition kinetic rate equation against
the high-overpressure experimental dataset of Kubo et al. (2002). This
regression constrains the activation energy and growth-rate prefactor
within the uncertainties reported in the main text.
Figure~\ref{fig:kubo-parameter-ranges} projects the \(\pm\) 2\(\sigma\)
uncertainty in \(\ln(K^{\prime})\) onto the kinetic prefactor
\(Z_\mathrm{ri}\) across grain sizes of 1--10 mm (1.7e-11 to 2.6e9
mol\(^2\) J\(^{-2}\) s\(^{-1}\)) and shows that both the \emph{Sluggish
Kinetics / Warm Slabs} (SW; 5.08e-11 to 5.08e-5 mol\(^2\) J\(^{-2}\)
s\(^{-1}\)) and \emph{Nominal Kinetics / Cold Slabs} (NC; 1.6e-3 to
1.6e3 mol\(^2\) J\(^{-2}\) s\(^{-1}\)) suites fall within the 95\%
confidence interval.

Figures \ref{fig:slab-scen04-set2} to \ref{fig:slab-scen11-set2} show
representative slab simulation snapshots spanning slow, intermediate,
and fast \(Z_\mathrm{ri}\) after 100 Ma across four kinematic and
rheological scenarios in each suite: scenarios 4, 0, 6, and 3 in the SW
suite (Figures \ref{fig:slab-scen04-set2} to \ref{fig:slab-scen03-set2})
and scenarios 12, 8, 14, and 11 in the NC suite (Figures
\ref{fig:slab-scen12-set2} to \ref{fig:slab-scen11-set2}). Together with
the reference snapshots for scenarios 4 and 12 in the main text (Figures
3 and 4), each scenario is presented in two field sets showing dynamic
temperature \(\hat{T}\), dynamic density \(\hat{\rho}\), and vertical
velocity \(\vec{u}_y\), and dynamic temperature \(\hat{T}\), log
viscosity \(\eta\), and post-spinel volume fraction \(X_\mathrm{ps}\).
These snapshots illustrate how the three kinetic regimes at the 660
(quasi-equilibrium, intermediate, and stagnation) interact with inflow
velocity, slab strength, and lower-mantle viscosity contrast across both
suites.

Table~\ref{tbl:depth-profile-summary} reports the kinetic prefactor
\(Z_\mathrm{ri}\), total 660 displacement (\(\delta z\)) and transition
width (\(\delta w\)), maximum reaction rate (\(\dot{X}_\mathrm{max}\)),
maximum descent velocity through the 660 (\(\vec{v}_\mathrm{max}\)),
thermal (\(\delta z_t\)) and kinetic (\(\delta z_k\)) displacement
contributions, and kinetic regime for all 132 slab simulations across
scenarios 0--19 (Table 1 in the main text). These values underlie the
quantitative trends and figures discussed in the main text.

\clearpage

\section*{Material Model (Figures S1--S4)}\label{sec:material-model}
\addcontentsline{toc}{section}{Material Model (Figures S1--S4)}

\begin{figure}
\centering
\includegraphics[width=1\linewidth,height=\textheight,keepaspectratio,alt={Reference
material properties used in ASPECT simulations showing density
(\textbackslash bar\{\textbackslash rho\}), thermal expansivity
(\textbackslash bar\{\textbackslash alpha\}), compressibility
(\textbackslash bar\{\textbackslash beta\}), and specific heat capacity
(\textbackslash bar\{C\}\_p). Colored profiles are for pure Mg olivine
(ol), wadsleyite (wd), ringwoodite (ri), and the post-spinel (ps)
assemblage bridgmanite + periclase. Profiles are computed using the
methods described in Section 2.2.1 in the main text. Gray bands indicate
the 410, 520, and 660
transitions.}]{../figs/adiabat/material-property-profiles.png}
\caption{Reference material properties used in ASPECT simulations
showing density (\(\bar{\rho}\)), thermal expansivity
(\(\bar{\alpha}\)), compressibility (\(\bar{\beta}\)), and specific heat
capacity (\(\bar{C}_p\)). Colored profiles are for pure Mg olivine (ol),
wadsleyite (wd), ringwoodite (ri), and the post-spinel (ps) assemblage
bridgmanite + periclase. Profiles are computed using the methods
described in Section 2.2.1 in the main text. Gray bands indicate the
410, 520, and 660 transitions.}\label{fig:material-property-profiles}
\end{figure}

\begin{figure}
\centering
\includegraphics[width=0.8\linewidth,height=\textheight,keepaspectratio,alt={Reference
thermodynamic properties used in our ASPECT simulations. Colored
profiles are for pure Mg olivine (ol), wadsleyite (wd), ringwoodite
(ri), and the post-spinel (ps) assemblage bridgmanite + periclase.
Profiles are computed using the same methods as described in Figure~
(see Section 2.2.1 in the main text). Gray bands indicate the 410, 520,
and 660
transitions.}]{../figs/adiabat/thermodynamic-property-profiles.png}
\caption{Reference thermodynamic properties used in our ASPECT
simulations. Colored profiles are for pure Mg olivine (ol), wadsleyite
(wd), ringwoodite (ri), and the post-spinel (ps) assemblage bridgmanite
+ periclase. Profiles are computed using the same methods as described
in Figure~\ref{fig:material-property-profiles} (see Section 2.2.1 in the
main text). Gray bands indicate the 410, 520, and 660
transitions.}\label{fig:thermodynamic-property-profiles}
\end{figure}

\begin{figure}
\centering
\includegraphics[width=0.68\linewidth,height=\textheight,keepaspectratio,alt={Monte
Carlo regression of high-overpressure experimental data from Kubo et al.
(2002). The best-fit activation energy E\^{}\{\textbackslash ast\} = 310
\textbackslash pm 233 kJ mol\^{}\{-1\} (2\textbackslash sigma) and
growth-rate prefactor \textbackslash ln(K\^{}\{\textbackslash prime\}) =
-29.8 \textbackslash pm 22.1 are both consistent with Kubo et al. (2002)
within uncertainty. The large uncertainties in E\^{}\{\textbackslash
ast\} and \textbackslash ln(K\^{}\{\textbackslash prime\}) reflect the
small number of data points and relatively large experimental T
uncertainties.}]{../figs/kubo/regression.png}
\caption{Monte Carlo regression of high-overpressure experimental data
from Kubo et al. (2002). The best-fit activation energy \(E^{\ast}\) =
310 \(\pm\) 233 kJ mol\(^{-1}\) (2\(\sigma\)) and growth-rate prefactor
\(\ln(K^{\prime})\) = -29.8 \(\pm\) 22.1 are both consistent with Kubo
et al. (2002) within uncertainty. The large uncertainties in
\(E^{\ast}\) and \(\ln(K^{\prime})\) reflect the small number of data
points and relatively large experimental \(T\)
uncertainties.}\label{fig:regression}
\end{figure}

\begin{figure}
\centering
\includegraphics[width=0.63\linewidth,height=\textheight,keepaspectratio,alt={Permitted
Z\_\textbackslash mathrm\{ri\} range across grain size d derived from
the Monte Carlo regression in Figure~ (Z\_\textbackslash mathrm\{ri\} =
\textbackslash text\{6.67\}\textbackslash, K\^{}\{\textbackslash
prime\}\textbackslash, /\textbackslash, d\textbackslash, \textbackslash
bar\{V\}\_\textbackslash mathrm\{ps\}\^{}2, with \textbackslash
ln(K\^{}\{\textbackslash prime\}) = -29.8 \textbackslash pm 22.1 at
2\textbackslash sigma and \textbackslash bar\{V\}\_\textbackslash
mathrm\{ps\} = 3.381e-5 m\^{}3 mol\^{}\{-1\}). The gray band and dashed
lines mark the 95\% confidence interval (1.7e-11 to 2.6e9 mol\^{}2
J\^{}\{-2\} s\^{}\{-1\} across d = 1--10 mm), and the solid black line
shows the best-fit trend (Z\_\textbackslash mathrm\{ri\} \textbackslash
sim 0.07--0.7 mol\^{}2 J\^{}\{-2\} s\^{}\{-1\}). Red and blue shaded
bands indicate the six-order-of-magnitude Z\_\textbackslash mathrm\{ri\}
ranges sampled in the SW (5.08e-11 to 5.08e-5 mol\^{}2 J\^{}\{-2\}
s\^{}\{-1\}) and NC (1.6e-3 to 1.6e3 mol\^{}2 J\^{}\{-2\} s\^{}\{-1\})
suites.}]{../figs/kubo/kubo-parameter-ranges.png}
\caption{Permitted \(Z_\mathrm{ri}\) range across grain size \(d\)
derived from the Monte Carlo regression in Figure~\ref{fig:regression}
(\(Z_\mathrm{ri} = \text{6.67}\, K^{\prime}\, /\, d\, \bar{V}_\mathrm{ps}^2\),
with \(\ln(K^{\prime})\) = -29.8 \(\pm\) 22.1 at 2\(\sigma\) and
\(\bar{V}_\mathrm{ps}\) = 3.381e-5 m\(^3\) mol\(^{-1}\)). The gray band
and dashed lines mark the 95\% confidence interval (1.7e-11 to 2.6e9
mol\(^2\) J\(^{-2}\) s\(^{-1}\) across \(d\) = 1--10 mm), and the solid
black line shows the best-fit trend (\(Z_\mathrm{ri} \sim\) 0.07--0.7
mol\(^2\) J\(^{-2}\) s\(^{-1}\)). Red and blue shaded bands indicate the
six-order-of-magnitude \(Z_\mathrm{ri}\) ranges sampled in the SW
(5.08e-11 to 5.08e-5 mol\(^2\) J\(^{-2}\) s\(^{-1}\)) and NC (1.6e-3 to
1.6e3 mol\(^2\) J\(^{-2}\) s\(^{-1}\))
suites.}\label{fig:kubo-parameter-ranges}
\end{figure}

\clearpage

\section*{Simulation Snapshots: Slabs (Figures
S5--S18)}\label{sec:simulation-snapshots}
\addcontentsline{toc}{section}{Simulation Snapshots: Slabs (Figures
S5--S18)}

\begin{figure}
\centering
\includegraphics[width=1\linewidth,height=\textheight,keepaspectratio,alt={SW
suite simulations for scenario 4 (\textbackslash Delta T\_0 = -500 K,
\textbackslash vec\{u\}\_\textbackslash mathrm\{in\} = 5 cm
yr\^{}\{-1\}, B = 5, b = 1; Table 1 in the main text) after 100 Ma,
showing dynamic temperature \textbackslash hat\{T\} (left column), log
viscosity \textbackslash eta (middle column), and post-spinel volume
fraction X\_\textbackslash mathrm\{ps\} (right column) across stagnation
(top row: Z\_\textbackslash mathrm\{ri\} = 5.08e-11 mol\^{}2 J\^{}\{-2\}
s\^{}\{-1\}), intermediate (middle row: Z\_\textbackslash mathrm\{ri\} =
1.90e-8 mol\^{}2 J\^{}\{-2\} s\^{}\{-1\}), and quasi-equilibrium (bottom
row: Z\_\textbackslash mathrm\{ri\} = 5.08e-5 mol\^{}2 J\^{}\{-2\}
s\^{}\{-1\}) kinetic regimes (companion to Figure 3 in the main text).
Thick black lines highlight the 10\% and 90\% wadsleyite and post-spinel
volume fraction contours used to define 410 and 660 displacement
(\textbackslash delta z) and width (\textbackslash delta w). The
background grid indicates ASPECT's adaptive mesh refinement near high
thermal gradients and phase transitions. The white trace indicates the
representative depth profile used to extract data, and the white arrows
mark where transition structure is
evaluated.}]{../figs/simulation/compositions/slab-set2-composition-scen04-0100.png}
\caption{SW suite simulations for scenario 4 (\(\Delta T_0\) = -500 K,
\(\vec{u}_\mathrm{in}\) = 5 cm yr\(^{-1}\), \(B\) = 5, \(b\) = 1; Table
1 in the main text) after 100 Ma, showing dynamic temperature
\(\hat{T}\) (left column), log viscosity \(\eta\) (middle column), and
post-spinel volume fraction \(X_\mathrm{ps}\) (right column) across
stagnation (top row: \(Z_\mathrm{ri}\) = 5.08e-11 mol\(^2\) J\(^{-2}\)
s\(^{-1}\)), intermediate (middle row: \(Z_\mathrm{ri}\) = 1.90e-8
mol\(^2\) J\(^{-2}\) s\(^{-1}\)), and quasi-equilibrium (bottom row:
\(Z_\mathrm{ri}\) = 5.08e-5 mol\(^2\) J\(^{-2}\) s\(^{-1}\)) kinetic
regimes (companion to Figure 3 in the main text). Thick black lines
highlight the 10\% and 90\% wadsleyite and post-spinel volume fraction
contours used to define 410 and 660 displacement (\(\delta z\)) and
width (\(\delta w\)). The background grid indicates ASPECT's adaptive
mesh refinement near high thermal gradients and phase transitions. The
white trace indicates the representative depth profile used to extract
data, and the white arrows mark where transition structure is
evaluated.}\label{fig:slab-scen04-set2}
\end{figure}

\clearpage

\begin{figure}
\centering
\includegraphics[width=1\linewidth,height=\textheight,keepaspectratio,alt={SW
suite simulations for scenario 0 (\textbackslash Delta T\_0 = -500 K,
\textbackslash vec\{u\}\_\textbackslash mathrm\{in\} = 5 cm
yr\^{}\{-1\}, B = 5, b = 50) after 100 Ma, showing dynamic temperature
\textbackslash hat\{T\} (left column), dynamic density \textbackslash
hat\{\textbackslash rho\} (middle column), and vertical velocity
\textbackslash vec\{u\}\_y (right column). All other annotations follow
Figure~. The 50-fold lower-mantle viscosity jump slows descent below the
0.5 cm yr\^{}\{-1\} stagnation threshold across all three
rows.}]{../figs/simulation/compositions/slab-set1-composition-scen00-0100.png}
\caption{SW suite simulations for scenario 0 (\(\Delta T_0\) = -500 K,
\(\vec{u}_\mathrm{in}\) = 5 cm yr\(^{-1}\), \(B\) = 5, \(b\) = 50) after
100 Ma, showing dynamic temperature \(\hat{T}\) (left column), dynamic
density \(\hat{\rho}\) (middle column), and vertical velocity
\(\vec{u}_y\) (right column). All other annotations follow
Figure~\ref{fig:slab-scen04-set2}. The 50-fold lower-mantle viscosity
jump slows descent below the 0.5 cm yr\(^{-1}\) stagnation threshold
across all three rows.}\label{fig:slab-scen00-set1}
\end{figure}

\clearpage

\begin{figure}
\centering
\includegraphics[width=1\linewidth,height=\textheight,keepaspectratio,alt={SW
suite simulations for scenario 0 (\textbackslash Delta T\_0 = -500 K,
\textbackslash vec\{u\}\_\textbackslash mathrm\{in\} = 5 cm
yr\^{}\{-1\}, B = 5, b = 50) after 100 Ma, showing \textbackslash
hat\{T\}, \textbackslash eta, and X\_\textbackslash mathrm\{ps\}
(companion to Figure~; annotations follow
Figure~).}]{../figs/simulation/compositions/slab-set2-composition-scen00-0100.png}
\caption{SW suite simulations for scenario 0 (\(\Delta T_0\) = -500 K,
\(\vec{u}_\mathrm{in}\) = 5 cm yr\(^{-1}\), \(B\) = 5, \(b\) = 50) after
100 Ma, showing \(\hat{T}\), \(\eta\), and \(X_\mathrm{ps}\) (companion
to Figure~\ref{fig:slab-scen00-set1}; annotations follow
Figure~\ref{fig:slab-scen04-set2}).}\label{fig:slab-scen00-set2}
\end{figure}

\clearpage

\begin{figure}
\centering
\includegraphics[width=1\linewidth,height=\textheight,keepaspectratio,alt={SW
suite simulations for scenario 6 (\textbackslash Delta T\_0 = -500 K,
\textbackslash vec\{u\}\_\textbackslash mathrm\{in\} = 10 cm
yr\^{}\{-1\}, B = 5, b = 1) after 100 Ma, showing \textbackslash
hat\{T\}, \textbackslash hat\{\textbackslash rho\}, and \textbackslash
vec\{u\}\_y (annotations follow Figure~). Faster inflow advects more
metastable ringwoodite to depth, broadening and deepening the 660 across
stagnation (top row), intermediate (middle row), and quasi-equilibrium
(bottom row)
regimes.}]{../figs/simulation/compositions/slab-set1-composition-scen06-0100.png}
\caption{SW suite simulations for scenario 6 (\(\Delta T_0\) = -500 K,
\(\vec{u}_\mathrm{in}\) = 10 cm yr\(^{-1}\), \(B\) = 5, \(b\) = 1) after
100 Ma, showing \(\hat{T}\), \(\hat{\rho}\), and \(\vec{u}_y\)
(annotations follow Figure~\ref{fig:slab-scen04-set2}). Faster inflow
advects more metastable ringwoodite to depth, broadening and deepening
the 660 across stagnation (top row), intermediate (middle row), and
quasi-equilibrium (bottom row) regimes.}\label{fig:slab-scen06-set1}
\end{figure}

\clearpage

\begin{figure}
\centering
\includegraphics[width=1\linewidth,height=\textheight,keepaspectratio,alt={SW
suite simulations for scenario 6 (\textbackslash Delta T\_0 = -500 K,
\textbackslash vec\{u\}\_\textbackslash mathrm\{in\} = 10 cm
yr\^{}\{-1\}, B = 5, b = 1) after 100 Ma, showing \textbackslash
hat\{T\}, \textbackslash eta, and X\_\textbackslash mathrm\{ps\}
(companion to Figure~; annotations follow
Figure~).}]{../figs/simulation/compositions/slab-set2-composition-scen06-0100.png}
\caption{SW suite simulations for scenario 6 (\(\Delta T_0\) = -500 K,
\(\vec{u}_\mathrm{in}\) = 10 cm yr\(^{-1}\), \(B\) = 5, \(b\) = 1) after
100 Ma, showing \(\hat{T}\), \(\eta\), and \(X_\mathrm{ps}\) (companion
to Figure~\ref{fig:slab-scen06-set1}; annotations follow
Figure~\ref{fig:slab-scen04-set2}).}\label{fig:slab-scen06-set2}
\end{figure}

\clearpage

\begin{figure}
\centering
\includegraphics[width=1\linewidth,height=\textheight,keepaspectratio,alt={SW
suite simulations for scenario 3 (\textbackslash Delta T\_0 = -500 K,
\textbackslash vec\{u\}\_\textbackslash mathrm\{in\} = 10 cm
yr\^{}\{-1\}, B = 10, b = 50) after 100 Ma, showing \textbackslash
hat\{T\}, \textbackslash hat\{\textbackslash rho\}, and \textbackslash
vec\{u\}\_y (annotations follow Figure~). High slab strength and fast
inflow produce the deepest (\textbackslash delta z = 80.5 km,
\textbackslash delta z\_k = 64.1 km) and widest (\textbackslash delta w
= 33.7 km) 660 in the stagnation regime (top row) while maintaining
penetration in the intermediate (middle row) and quasi-equilibrium
(bottom row) regimes despite b =
50.}]{../figs/simulation/compositions/slab-set1-composition-scen03-0100.png}
\caption{SW suite simulations for scenario 3 (\(\Delta T_0\) = -500 K,
\(\vec{u}_\mathrm{in}\) = 10 cm yr\(^{-1}\), \(B\) = 10, \(b\) = 50)
after 100 Ma, showing \(\hat{T}\), \(\hat{\rho}\), and \(\vec{u}_y\)
(annotations follow Figure~\ref{fig:slab-scen04-set2}). High slab
strength and fast inflow produce the deepest (\(\delta z\) = 80.5 km,
\(\delta z_k\) = 64.1 km) and widest (\(\delta w\) = 33.7 km) 660 in the
stagnation regime (top row) while maintaining penetration in the
intermediate (middle row) and quasi-equilibrium (bottom row) regimes
despite \(b\) = 50.}\label{fig:slab-scen03-set1}
\end{figure}

\clearpage

\begin{figure}
\centering
\includegraphics[width=1\linewidth,height=\textheight,keepaspectratio,alt={SW
suite simulations for scenario 3 (\textbackslash Delta T\_0 = -500 K,
\textbackslash vec\{u\}\_\textbackslash mathrm\{in\} = 10 cm
yr\^{}\{-1\}, B = 10, b = 50) after 100 Ma, showing \textbackslash
hat\{T\}, \textbackslash eta, and X\_\textbackslash mathrm\{ps\}
(companion to Figure~; annotations follow
Figure~).}]{../figs/simulation/compositions/slab-set2-composition-scen03-0100.png}
\caption{SW suite simulations for scenario 3 (\(\Delta T_0\) = -500 K,
\(\vec{u}_\mathrm{in}\) = 10 cm yr\(^{-1}\), \(B\) = 10, \(b\) = 50)
after 100 Ma, showing \(\hat{T}\), \(\eta\), and \(X_\mathrm{ps}\)
(companion to Figure~\ref{fig:slab-scen03-set1}; annotations follow
Figure~\ref{fig:slab-scen04-set2}).}\label{fig:slab-scen03-set2}
\end{figure}

\clearpage

\begin{figure}
\centering
\includegraphics[width=1\linewidth,height=\textheight,keepaspectratio,alt={NC
suite simulations for scenario 12 (\textbackslash Delta T\_0 = -1000 K,
\textbackslash vec\{u\}\_\textbackslash mathrm\{in\} = 5 cm
yr\^{}\{-1\}, B = 5, b = 1) after 100 Ma, showing \textbackslash
hat\{T\}, \textbackslash eta, and X\_\textbackslash mathrm\{ps\} across
Z\_\textbackslash mathrm\{ri\} = 1.6e-3 (top row), 6.0e-1 (middle row),
and 1.6e3 (bottom row) mol\^{}2 J\^{}\{-2\} s\^{}\{-1\} (companion to
Figure 4 in the main text; annotations follow
Figure~).}]{../figs/simulation/compositions/slab-set2-composition-scen12-0100.png}
\caption{NC suite simulations for scenario 12 (\(\Delta T_0\) = -1000 K,
\(\vec{u}_\mathrm{in}\) = 5 cm yr\(^{-1}\), \(B\) = 5, \(b\) = 1) after
100 Ma, showing \(\hat{T}\), \(\eta\), and \(X_\mathrm{ps}\) across
\(Z_\mathrm{ri}\) = 1.6e-3 (top row), 6.0e-1 (middle row), and 1.6e3
(bottom row) mol\(^2\) J\(^{-2}\) s\(^{-1}\) (companion to Figure 4 in
the main text; annotations follow
Figure~\ref{fig:slab-scen04-set2}).}\label{fig:slab-scen12-set2}
\end{figure}

\clearpage

\begin{figure}
\centering
\includegraphics[width=1\linewidth,height=\textheight,keepaspectratio,alt={NC
suite simulations for scenario 8 (\textbackslash Delta T\_0 = -1000 K,
\textbackslash vec\{u\}\_\textbackslash mathrm\{in\} = 5 cm
yr\^{}\{-1\}, B = 5, b = 50) after 100 Ma, showing \textbackslash
hat\{T\}, \textbackslash hat\{\textbackslash rho\}, and \textbackslash
vec\{u\}\_y (annotations follow Figure~). The 50-fold lower-mantle
viscosity jump slows and folds the cold slab without inducing kinetic
broadening or stagnation across any tested Z\_\textbackslash
mathrm\{ri\}.}]{../figs/simulation/compositions/slab-set1-composition-scen08-0100.png}
\caption{NC suite simulations for scenario 8 (\(\Delta T_0\) = -1000 K,
\(\vec{u}_\mathrm{in}\) = 5 cm yr\(^{-1}\), \(B\) = 5, \(b\) = 50) after
100 Ma, showing \(\hat{T}\), \(\hat{\rho}\), and \(\vec{u}_y\)
(annotations follow Figure~\ref{fig:slab-scen04-set2}). The 50-fold
lower-mantle viscosity jump slows and folds the cold slab without
inducing kinetic broadening or stagnation across any tested
\(Z_\mathrm{ri}\).}\label{fig:slab-scen08-set1}
\end{figure}

\clearpage

\begin{figure}
\centering
\includegraphics[width=1\linewidth,height=\textheight,keepaspectratio,alt={NC
suite simulations for scenario 8 (\textbackslash Delta T\_0 = -1000 K,
\textbackslash vec\{u\}\_\textbackslash mathrm\{in\} = 5 cm
yr\^{}\{-1\}, B = 5, b = 50) after 100 Ma, showing \textbackslash
hat\{T\}, \textbackslash eta, and X\_\textbackslash mathrm\{ps\}
(companion to Figure~; annotations follow
Figure~).}]{../figs/simulation/compositions/slab-set2-composition-scen08-0100.png}
\caption{NC suite simulations for scenario 8 (\(\Delta T_0\) = -1000 K,
\(\vec{u}_\mathrm{in}\) = 5 cm yr\(^{-1}\), \(B\) = 5, \(b\) = 50) after
100 Ma, showing \(\hat{T}\), \(\eta\), and \(X_\mathrm{ps}\) (companion
to Figure~\ref{fig:slab-scen08-set1}; annotations follow
Figure~\ref{fig:slab-scen04-set2}).}\label{fig:slab-scen08-set2}
\end{figure}

\clearpage

\begin{figure}
\centering
\includegraphics[width=1\linewidth,height=\textheight,keepaspectratio,alt={NC
suite simulations for scenario 14 (\textbackslash Delta T\_0 = -1000 K,
\textbackslash vec\{u\}\_\textbackslash mathrm\{in\} = 10 cm
yr\^{}\{-1\}, B = 5, b = 1) after 100 Ma, showing \textbackslash
hat\{T\}, \textbackslash hat\{\textbackslash rho\}, and \textbackslash
vec\{u\}\_y (annotations follow Figure~). Fast descent (\textbackslash
vec\{v\}\_\textbackslash mathrm\{max\} = 4.8--6.6 cm yr\^{}\{-1\})
drives the slowest nominal Z\_\textbackslash mathrm\{ri\} (top row:
Z\_\textbackslash mathrm\{ri\} = 1.6e-3 mol\^{}2 J\^{}\{-2\}
s\^{}\{-1\}) into the intermediate regime (L = 13.0 km, \textbackslash
delta z\_k = 22.3 km, \textbackslash delta w = 13.1 km), while faster
Z\_\textbackslash mathrm\{ri\} rows (middle and bottom) remain in
quasi-equilibrium.}]{../figs/simulation/compositions/slab-set1-composition-scen14-0100.png}
\caption{NC suite simulations for scenario 14 (\(\Delta T_0\) = -1000 K,
\(\vec{u}_\mathrm{in}\) = 10 cm yr\(^{-1}\), \(B\) = 5, \(b\) = 1) after
100 Ma, showing \(\hat{T}\), \(\hat{\rho}\), and \(\vec{u}_y\)
(annotations follow Figure~\ref{fig:slab-scen04-set2}). Fast descent
(\(\vec{v}_\mathrm{max}\) = 4.8--6.6 cm yr\(^{-1}\)) drives the slowest
nominal \(Z_\mathrm{ri}\) (top row: \(Z_\mathrm{ri}\) = 1.6e-3 mol\(^2\)
J\(^{-2}\) s\(^{-1}\)) into the intermediate regime (\(L\) = 13.0 km,
\(\delta z_k\) = 22.3 km, \(\delta w\) = 13.1 km), while faster
\(Z_\mathrm{ri}\) rows (middle and bottom) remain in
quasi-equilibrium.}\label{fig:slab-scen14-set1}
\end{figure}

\clearpage

\begin{figure}
\centering
\includegraphics[width=1\linewidth,height=\textheight,keepaspectratio,alt={NC
suite simulations for scenario 14 (\textbackslash Delta T\_0 = -1000 K,
\textbackslash vec\{u\}\_\textbackslash mathrm\{in\} = 10 cm
yr\^{}\{-1\}, B = 5, b = 1) after 100 Ma, showing \textbackslash
hat\{T\}, \textbackslash eta, and X\_\textbackslash mathrm\{ps\}
(companion to Figure~; annotations follow
Figure~).}]{../figs/simulation/compositions/slab-set2-composition-scen14-0100.png}
\caption{NC suite simulations for scenario 14 (\(\Delta T_0\) = -1000 K,
\(\vec{u}_\mathrm{in}\) = 10 cm yr\(^{-1}\), \(B\) = 5, \(b\) = 1) after
100 Ma, showing \(\hat{T}\), \(\eta\), and \(X_\mathrm{ps}\) (companion
to Figure~\ref{fig:slab-scen14-set1}; annotations follow
Figure~\ref{fig:slab-scen04-set2}).}\label{fig:slab-scen14-set2}
\end{figure}

\clearpage

\begin{figure}
\centering
\includegraphics[width=1\linewidth,height=\textheight,keepaspectratio,alt={NC
suite simulations for scenario 11 (\textbackslash Delta T\_0 = -1000 K,
\textbackslash vec\{u\}\_\textbackslash mathrm\{in\} = 10 cm
yr\^{}\{-1\}, B = 10, b = 50) after 100 Ma, showing \textbackslash
hat\{T\}, \textbackslash hat\{\textbackslash rho\}, and \textbackslash
vec\{u\}\_y (annotations follow Figure~). All three rows remain in the
quasi-equilibrium regime with \textbackslash delta z\_k = 5.8--7.3
km.}]{../figs/simulation/compositions/slab-set1-composition-scen11-0100.png}
\caption{NC suite simulations for scenario 11 (\(\Delta T_0\) = -1000 K,
\(\vec{u}_\mathrm{in}\) = 10 cm yr\(^{-1}\), \(B\) = 10, \(b\) = 50)
after 100 Ma, showing \(\hat{T}\), \(\hat{\rho}\), and \(\vec{u}_y\)
(annotations follow Figure~\ref{fig:slab-scen04-set2}). All three rows
remain in the quasi-equilibrium regime with \(\delta z_k\) = 5.8--7.3
km.}\label{fig:slab-scen11-set1}
\end{figure}

\clearpage

\begin{figure}
\centering
\includegraphics[width=1\linewidth,height=\textheight,keepaspectratio,alt={NC
suite simulations for scenario 11 (\textbackslash Delta T\_0 = -1000 K,
\textbackslash vec\{u\}\_\textbackslash mathrm\{in\} = 10 cm
yr\^{}\{-1\}, B = 10, b = 50) after 100 Ma, showing \textbackslash
hat\{T\}, \textbackslash eta, and X\_\textbackslash mathrm\{ps\}
(companion to Figure~; annotations follow
Figure~).}]{../figs/simulation/compositions/slab-set2-composition-scen11-0100.png}
\caption{NC suite simulations for scenario 11 (\(\Delta T_0\) = -1000 K,
\(\vec{u}_\mathrm{in}\) = 10 cm yr\(^{-1}\), \(B\) = 10, \(b\) = 50)
after 100 Ma, showing \(\hat{T}\), \(\eta\), and \(X_\mathrm{ps}\)
(companion to Figure~\ref{fig:slab-scen11-set1}; annotations follow
Figure~\ref{fig:slab-scen04-set2}).}\label{fig:slab-scen11-set2}
\end{figure}

\clearpage

\section*{Structure of the 660: Displacement, Width, and Thermal-Kinetic
Contributions}\label{sec:structure-of-the-660}
\addcontentsline{toc}{section}{Structure of the 660: Displacement,
Width, and Thermal-Kinetic Contributions}

\begin{longtable}[]{@{}
  >{\raggedright\arraybackslash}p{(\linewidth - 18\tabcolsep) * \real{0.0448}}
  >{\raggedleft\arraybackslash}p{(\linewidth - 18\tabcolsep) * \real{0.0672}}
  >{\raggedleft\arraybackslash}p{(\linewidth - 18\tabcolsep) * \real{0.1194}}
  >{\raggedleft\arraybackslash}p{(\linewidth - 18\tabcolsep) * \real{0.0970}}
  >{\raggedleft\arraybackslash}p{(\linewidth - 18\tabcolsep) * \real{0.0821}}
  >{\raggedleft\arraybackslash}p{(\linewidth - 18\tabcolsep) * \real{0.1716}}
  >{\raggedleft\arraybackslash}p{(\linewidth - 18\tabcolsep) * \real{0.1716}}
  >{\raggedleft\arraybackslash}p{(\linewidth - 18\tabcolsep) * \real{0.0970}}
  >{\raggedleft\arraybackslash}p{(\linewidth - 18\tabcolsep) * \real{0.0970}}
  >{\raggedleft\arraybackslash}p{(\linewidth - 18\tabcolsep) * \real{0.0522}}@{}}
\caption{Summary of kinetic prefactor \(Z_\mathrm{ri}\), total 660
displacement (\(\delta z\)) and transition width (\(\delta w\)), maximum
reaction rate (\(\dot{X}_\mathrm{max}\)), maximum descent velocity
through the 660 (\(\vec{v}_\mathrm{max}\)), thermal (\(\delta z_t\)) and
kinetic (\(\delta z_k\)) displacement contributions, and kinetic regime
across all 20 scenarios (SW: \emph{Sluggish Kinetics / Warm Slabs},
scenarios 0--7; NC: \emph{Nominal Kinetics / Cold Slabs}, scenarios
8--15; T1--T4: sluggish-410 sensitivity tests, scenarios 16--19; see
Table 1 in the main text). Units are \(Z_\mathrm{ri}\): mol\(^2\)
J\(^{-2}\) s\(^{-1}\); \(\delta z\), \(\delta w\), \(\delta z_t\), and
\(\delta z_k\): km; \(\dot{X}_\mathrm{max}\): Ma\(^{-1}\);
\(\vec{v}_\mathrm{max}\): cm yr\(^{-1}\). Regimes correspond to 1:
quasi-equilibrium, 2: intermediate, and 3:
stagnation.}\label{tbl:depth-profile-summary}\tabularnewline
\toprule\noalign{}
\begin{minipage}[b]{\linewidth}\raggedright
Suite
\end{minipage} & \begin{minipage}[b]{\linewidth}\raggedleft
Scenario
\end{minipage} & \begin{minipage}[b]{\linewidth}\raggedleft
\(Z_\mathrm{ri}\)
\end{minipage} & \begin{minipage}[b]{\linewidth}\raggedleft
\(\delta z\)
\end{minipage} & \begin{minipage}[b]{\linewidth}\raggedleft
\(\delta w\)
\end{minipage} & \begin{minipage}[b]{\linewidth}\raggedleft
\(\dot{X}_\mathrm{max}\)
\end{minipage} & \begin{minipage}[b]{\linewidth}\raggedleft
\(\vec{v}_\mathrm{max}\)
\end{minipage} & \begin{minipage}[b]{\linewidth}\raggedleft
\(\delta z_t\)
\end{minipage} & \begin{minipage}[b]{\linewidth}\raggedleft
\(\delta z_k\)
\end{minipage} & \begin{minipage}[b]{\linewidth}\raggedleft
Regime
\end{minipage} \\
\midrule\noalign{}
\endfirsthead
\toprule\noalign{}
\begin{minipage}[b]{\linewidth}\raggedright
Suite
\end{minipage} & \begin{minipage}[b]{\linewidth}\raggedleft
Scenario
\end{minipage} & \begin{minipage}[b]{\linewidth}\raggedleft
\(Z_\mathrm{ri}\)
\end{minipage} & \begin{minipage}[b]{\linewidth}\raggedleft
\(\delta z\)
\end{minipage} & \begin{minipage}[b]{\linewidth}\raggedleft
\(\delta w\)
\end{minipage} & \begin{minipage}[b]{\linewidth}\raggedleft
\(\dot{X}_\mathrm{max}\)
\end{minipage} & \begin{minipage}[b]{\linewidth}\raggedleft
\(\vec{v}_\mathrm{max}\)
\end{minipage} & \begin{minipage}[b]{\linewidth}\raggedleft
\(\delta z_t\)
\end{minipage} & \begin{minipage}[b]{\linewidth}\raggedleft
\(\delta z_k\)
\end{minipage} & \begin{minipage}[b]{\linewidth}\raggedleft
Regime
\end{minipage} \\
\midrule\noalign{}
\endhead
\bottomrule\noalign{}
\endlastfoot
SW & 0 & 5.1e-05 & 26.7 & 11.0 & 1.35e+02 & 0.49 & 18.0 & 8.7 & 3 \\
SW & 0 & 7.0e-06 & 27.3 & 11.5 & 1.75e+01 & 0.48 & 18.0 & 9.3 & 3 \\
SW & 0 & 9.8e-07 & 29.0 & 6.5 & 1.67e+00 & 0.32 & 17.9 & 11.1 & 3 \\
SW & 0 & 1.4e-07 & 31.1 & 6.1 & 6.85e-01 & 0.41 & 17.7 & 13.3 & 3 \\
SW & 0 & 1.9e-08 & 35.2 & 8.3 & 3.22e-01 & 0.30 & 17.7 & 17.5 & 3 \\
SW & 0 & 2.6e-09 & 40.7 & 11.8 & 1.58e-01 & 0.20 & 17.1 & 23.6 & 3 \\
SW & 0 & 3.8e-10 & 46.6 & 15.0 & 7.70e-02 & 0.14 & 15.5 & 31.1 & 3 \\
SW & 0 & 5.1e-11 & 51.4 & 16.9 & 4.12e-02 & 0.09 & 12.8 & 38.6 & 3 \\
SW & 1 & 5.1e-05 & 22.5 & 4.0 & 7.16e+00 & 0.61 & 15.7 & 6.7 & 1 \\
SW & 1 & 7.0e-06 & 27.1 & 11.3 & 1.66e+01 & 0.58 & 16.9 & 10.1 & 1 \\
SW & 1 & 9.8e-07 & 28.2 & 12.2 & 2.47e+00 & 0.56 & 16.9 & 11.3 & 1 \\
SW & 1 & 1.4e-07 & 30.1 & 7.1 & 1.24e+00 & 0.54 & 16.5 & 13.6 & 1 \\
SW & 1 & 1.9e-08 & 34.6 & 7.8 & 5.03e-01 & 0.46 & 16.8 & 17.8 & 3 \\
SW & 1 & 2.6e-09 & 42.9 & 13.6 & 2.24e-01 & 0.32 & 16.8 & 26.1 & 3 \\
SW & 1 & 3.8e-10 & 50.8 & 18.0 & 9.75e-02 & 0.20 & 15.9 & 35.0 & 3 \\
SW & 1 & 5.1e-11 & 58.7 & 21.4 & 5.00e-02 & 0.13 & 13.6 & 45.1 & 3 \\
SW & 2 & 5.1e-05 & 30.0 & 5.3 & 1.91e+01 & 0.94 & 20.3 & 9.7 & 1 \\
SW & 2 & 7.0e-06 & 30.4 & 4.6 & 3.15e+00 & 0.91 & 20.2 & 10.2 & 1 \\
SW & 2 & 9.8e-07 & 33.4 & 6.7 & 1.29e+00 & 0.87 & 20.9 & 12.4 & 2 \\
SW & 2 & 1.4e-07 & 39.7 & 9.4 & 9.06e-01 & 0.80 & 21.1 & 18.6 & 2 \\
SW & 2 & 1.9e-08 & 48.9 & 14.7 & 3.95e-01 & 0.61 & 21.2 & 27.7 & 2 \\
SW & 2 & 2.6e-09 & 60.9 & 22.0 & 1.52e-01 & 0.37 & 21.1 & 39.8 & 3 \\
SW & 2 & 3.8e-10 & 68.9 & 25.3 & 6.15e-02 & 0.20 & 19.6 & 49.2 & 3 \\
SW & 2 & 5.1e-11 & 74.9 & 24.7 & 3.81e-02 & 0.14 & 16.4 & 58.5 & 3 \\
SW & 3 & 5.1e-05 & 22.4 & 4.0 & 3.37e+01 & 1.07 & 14.6 & 7.9 & 1 \\
SW & 3 & 7.0e-06 & 22.5 & 4.0 & 6.08e+00 & 1.10 & 14.5 & 8.0 & 1 \\
SW & 3 & 9.8e-07 & 28.2 & 12.1 & 3.59e+00 & 1.05 & 16.3 & 11.9 & 1 \\
SW & 3 & 1.4e-07 & 30.4 & 6.8 & 2.00e+00 & 1.02 & 15.7 & 14.7 & 2 \\
SW & 3 & 1.9e-08 & 39.8 & 11.6 & 9.13e-01 & 0.99 & 17.2 & 22.7 & 2 \\
SW & 3 & 2.6e-09 & 52.9 & 19.1 & 3.38e-01 & 0.67 & 17.9 & 35.0 & 2 \\
SW & 3 & 3.8e-10 & 65.6 & 26.7 & 1.24e-01 & 0.36 & 17.7 & 47.9 & 3 \\
SW & 3 & 5.1e-11 & 80.5 & 33.7 & 5.13e-02 & 0.22 & 16.4 & 64.1 & 3 \\
SW & 4 & 5.1e-05 & 25.1 & 6.3 & 3.44e+00 & 1.55 & 16.5 & 8.6 & 1 \\
SW & 4 & 7.0e-06 & 28.2 & 6.0 & 2.14e+01 & 1.73 & 16.7 & 11.5 & 1 \\
SW & 4 & 9.8e-07 & 29.9 & 7.2 & 6.39e+00 & 1.87 & 16.9 & 13.1 & 1 \\
SW & 4 & 1.4e-07 & 33.1 & 7.0 & 2.24e+00 & 1.68 & 17.0 & 16.1 & 2 \\
SW & 4 & 1.9e-08 & 39.6 & 11.8 & 9.12e-01 & 1.11 & 17.1 & 22.4 & 2 \\
SW & 4 & 2.6e-09 & 44.0 & 13.4 & 2.17e-01 & 0.49 & 17.2 & 26.9 & 3 \\
SW & 4 & 3.8e-10 & 43.1 & 10.7 & 1.12e-01 & 0.31 & 13.4 & 29.7 & 3 \\
SW & 4 & 5.1e-11 & 48.4 & 10.8 & 6.82e-02 & 0.25 & 11.0 & 37.4 & 3 \\
SW & 5 & 5.1e-05 & 24.4 & 5.8 & 6.03e+01 & 2.05 & 15.0 & 9.5 & 1 \\
SW & 5 & 7.0e-06 & 24.7 & 6.0 & 7.22e+00 & 1.92 & 15.2 & 9.5 & 1 \\
SW & 5 & 9.8e-07 & 27.6 & 6.8 & 4.52e+00 & 1.73 & 15.7 & 11.9 & 1 \\
SW & 5 & 1.4e-07 & 30.0 & 7.1 & 2.22e+00 & 1.17 & 15.6 & 14.4 & 2 \\
SW & 5 & 1.9e-08 & 39.6 & 12.2 & 1.41e+00 & 1.71 & 16.2 & 23.4 & 2 \\
SW & 5 & 2.6e-09 & 52.4 & 20.0 & 6.38e-01 & 1.31 & 16.3 & 36.1 & 2 \\
SW & 5 & 3.8e-10 & 55.8 & 20.1 & 1.36e-01 & 0.46 & 16.0 & 39.8 & 3 \\
SW & 5 & 5.1e-11 & 57.8 & 16.4 & 7.96e-02 & 0.29 & 12.1 & 45.7 & 3 \\
SW & 6 & 5.1e-05 & 31.4 & 5.4 & 2.83e+01 & 3.21 & 20.6 & 10.7 & 1 \\
SW & 6 & 7.0e-06 & 32.5 & 6.2 & 7.40e+00 & 3.07 & 20.7 & 11.9 & 1 \\
SW & 6 & 9.8e-07 & 36.5 & 8.9 & 3.45e+00 & 2.85 & 20.7 & 15.8 & 2 \\
SW & 6 & 1.4e-07 & 44.0 & 12.2 & 1.84e+00 & 2.39 & 20.9 & 23.2 & 2 \\
SW & 6 & 1.9e-08 & 54.7 & 18.7 & 7.04e-01 & 1.44 & 21.1 & 33.7 & 2 \\
SW & 6 & 2.6e-09 & 55.7 & 17.8 & 1.36e-01 & 0.44 & 20.6 & 35.1 & 3 \\
SW & 6 & 3.8e-10 & 54.5 & 13.1 & 1.06e-01 & 0.36 & 15.9 & 38.5 & 3 \\
SW & 6 & 5.1e-11 & 68.4 & 13.4 & 7.39e-02 & 0.32 & 13.9 & 54.5 & 3 \\
SW & 7 & 5.1e-05 & 28.1 & 12.0 & 1.04e+02 & 2.29 & 17.4 & 10.7 & 1 \\
SW & 7 & 7.0e-06 & 28.4 & 12.3 & 1.38e+01 & 2.35 & 17.4 & 11.1 & 1 \\
SW & 7 & 9.8e-07 & 30.0 & 7.1 & 5.71e+00 & 2.19 & 17.1 & 12.9 & 1 \\
SW & 7 & 1.4e-07 & 34.7 & 7.9 & 2.25e+00 & 2.01 & 17.7 & 17.0 & 2 \\
SW & 7 & 1.9e-08 & 44.3 & 14.6 & 1.37e+00 & 2.09 & 17.5 & 26.8 & 2 \\
SW & 7 & 2.6e-09 & 60.3 & 23.8 & 6.14e-01 & 1.59 & 17.7 & 42.6 & 2 \\
SW & 7 & 3.8e-10 & 79.6 & 34.3 & 1.94e-01 & 0.91 & 18.3 & 61.4 & 2 \\
SW & 7 & 5.1e-11 & 78.3 & 26.1 & 7.63e-02 & 0.48 & 14.9 & 63.5 & 3 \\
NC & 8 & 1.6e+03 & 35.8 & 5.9 & 9.49e+03 & 0.98 & 30.2 & 5.6 & 1 \\
NC & 8 & 2.2e+02 & 35.9 & 6.0 & 6.82e+03 & 0.99 & 29.9 & 6.0 & 1 \\
NC & 8 & 3.1e+01 & 39.0 & 5.4 & 9.14e+04 & 1.04 & 30.7 & 8.4 & 1 \\
NC & 8 & 4.3e+00 & 39.0 & 5.4 & 1.11e+04 & 1.01 & 30.8 & 8.2 & 1 \\
NC & 8 & 6.0e-01 & 39.0 & 5.4 & 1.67e+03 & 1.01 & 30.7 & 8.3 & 1 \\
NC & 8 & 8.3e-02 & 39.0 & 5.4 & 2.35e+02 & 1.01 & 30.7 & 8.3 & 1 \\
NC & 8 & 1.2e-02 & 39.0 & 5.4 & 3.41e+01 & 1.02 & 30.7 & 8.3 & 1 \\
NC & 8 & 1.6e-03 & 39.3 & 5.7 & 4.72e+00 & 1.02 & 30.8 & 8.6 & 1 \\
NC & 9 & 1.6e+03 & 27.9 & 12.0 & 2.27e+09 & 1.03 & 20.8 & 7.1 & 1 \\
NC & 9 & 2.2e+02 & 28.1 & 12.2 & 3.18e+08 & 0.98 & 21.0 & 7.1 & 1 \\
NC & 9 & 3.1e+01 & 27.9 & 12.0 & 4.53e+07 & 0.96 & 21.0 & 7.0 & 1 \\
NC & 9 & 4.3e+00 & 29.8 & 13.8 & 5.79e+06 & 0.97 & 20.8 & 9.0 & 1 \\
NC & 9 & 6.0e-01 & 27.9 & 12.0 & 8.81e+05 & 0.93 & 20.9 & 7.0 & 1 \\
NC & 9 & 8.3e-02 & 27.9 & 11.9 & 1.11e+05 & 0.96 & 21.0 & 7.0 & 1 \\
NC & 9 & 1.2e-02 & 29.3 & 13.2 & 1.51e+04 & 1.00 & 20.8 & 8.4 & 1 \\
NC & 9 & 1.6e-03 & 29.7 & 4.8 & 9.72e+02 & 1.09 & 20.5 & 9.2 & 1 \\
NC & 10 & 1.6e+03 & 36.4 & 3.2 & 1.21e+07 & 1.55 & 29.6 & 6.8 & 1 \\
NC & 10 & 2.2e+02 & 36.4 & 3.2 & 1.51e+06 & 1.51 & 29.1 & 7.3 & 1 \\
NC & 10 & 3.1e+01 & 38.6 & 5.0 & 1.53e+04 & 1.69 & 32.0 & 6.5 & 1 \\
NC & 10 & 4.3e+00 & 38.6 & 5.0 & 1.83e+03 & 1.67 & 32.1 & 6.4 & 1 \\
NC & 10 & 6.0e-01 & 38.6 & 5.0 & 2.39e+02 & 1.73 & 32.2 & 6.4 & 1 \\
NC & 10 & 8.3e-02 & 38.8 & 5.2 & 7.66e+00 & 1.61 & 33.2 & 5.6 & 1 \\
NC & 10 & 1.2e-02 & 45.5 & 6.7 & 4.57e+00 & 1.49 & 35.2 & 10.2 & 1 \\
NC & 10 & 1.6e-03 & 46.2 & 5.3 & 3.99e+00 & 1.56 & 33.7 & 12.5 & 1 \\
NC & 11 & 1.6e+03 & 22.6 & 4.1 & 7.27e+06 & 1.50 & 16.7 & 5.8 & 1 \\
NC & 11 & 2.2e+02 & 22.8 & 4.3 & 2.00e+05 & 1.67 & 16.7 & 6.2 & 1 \\
NC & 11 & 3.1e+01 & 22.8 & 4.3 & 5.46e+04 & 1.55 & 16.7 & 6.2 & 1 \\
NC & 11 & 4.3e+00 & 22.7 & 4.2 & 2.02e+04 & 1.59 & 16.6 & 6.1 & 1 \\
NC & 11 & 6.0e-01 & 22.5 & 4.0 & 1.15e+03 & 1.51 & 16.7 & 5.8 & 1 \\
NC & 11 & 8.3e-02 & 22.8 & 4.3 & 2.62e+02 & 1.53 & 16.6 & 6.2 & 1 \\
NC & 11 & 1.2e-02 & 27.8 & 11.9 & 1.86e+04 & 1.85 & 20.4 & 7.3 & 1 \\
NC & 11 & 1.6e-03 & 27.6 & 11.4 & 2.07e+03 & 1.56 & 20.3 & 7.3 & 1 \\
NC & 12 & 1.6e+03 & 37.4 & 8.1 & 3.26e+07 & 3.04 & 29.6 & 7.8 & 1 \\
NC & 12 & 2.2e+02 & 37.4 & 8.1 & 4.49e+06 & 3.04 & 29.6 & 7.8 & 1 \\
NC & 12 & 3.1e+01 & 37.4 & 8.1 & 6.32e+05 & 3.04 & 29.6 & 7.8 & 1 \\
NC & 12 & 4.3e+00 & 37.4 & 8.1 & 8.78e+04 & 3.06 & 29.6 & 7.8 & 1 \\
NC & 12 & 6.0e-01 & 37.4 & 8.1 & 1.22e+04 & 3.05 & 29.6 & 7.8 & 1 \\
NC & 12 & 8.3e-02 & 37.4 & 8.1 & 1.69e+03 & 3.02 & 29.6 & 7.8 & 1 \\
NC & 12 & 1.2e-02 & 38.4 & 7.5 & 1.93e+02 & 3.04 & 29.6 & 8.9 & 1 \\
NC & 12 & 1.6e-03 & 40.1 & 6.3 & 2.44e+01 & 2.96 & 29.6 & 10.5 & 1 \\
NC & 13 & 1.6e+03 & 32.6 & 6.2 & 6.54e+08 & 2.93 & 21.8 & 10.8 & 1 \\
NC & 13 & 2.2e+02 & 32.6 & 6.2 & 9.05e+07 & 2.93 & 21.8 & 10.8 & 1 \\
NC & 13 & 3.1e+01 & 32.6 & 6.2 & 1.26e+07 & 2.91 & 21.8 & 10.8 & 1 \\
NC & 13 & 4.3e+00 & 32.6 & 6.2 & 1.77e+06 & 2.93 & 21.8 & 10.8 & 1 \\
NC & 13 & 6.0e-01 & 32.6 & 6.2 & 2.46e+05 & 2.92 & 21.8 & 10.8 & 1 \\
NC & 13 & 8.3e-02 & 32.6 & 6.2 & 3.38e+04 & 2.92 & 21.8 & 10.8 & 1 \\
NC & 13 & 1.2e-02 & 32.6 & 6.2 & 4.85e+03 & 2.94 & 21.8 & 10.8 & 1 \\
NC & 13 & 1.6e-03 & 32.6 & 6.2 & 6.38e+02 & 2.92 & 21.8 & 10.8 & 1 \\
NC & 14 & 1.6e+03 & 45.7 & 8.3 & 9.47e+04 & 6.34 & 38.8 & 6.9 & 1 \\
NC & 14 & 2.2e+02 & 45.7 & 8.3 & 1.33e+04 & 6.31 & 38.8 & 6.9 & 1 \\
NC & 14 & 3.1e+01 & 45.9 & 8.3 & 1.64e+03 & 6.62 & 39.0 & 6.9 & 1 \\
NC & 14 & 4.3e+00 & 46.6 & 7.6 & 1.86e+02 & 6.39 & 38.8 & 7.8 & 1 \\
NC & 14 & 6.0e-01 & 46.7 & 7.3 & 4.47e+01 & 5.53 & 38.1 & 8.6 & 1 \\
NC & 14 & 8.3e-02 & 48.1 & 6.5 & 1.46e+01 & 5.37 & 38.0 & 10.1 & 1 \\
NC & 14 & 1.2e-02 & 53.8 & 10.2 & 7.02e+00 & 5.52 & 38.3 & 15.5 & 2 \\
NC & 14 & 1.6e-03 & 60.1 & 13.1 & 3.72e+00 & 4.83 & 37.8 & 22.3 & 2 \\
NC & 15 & 1.6e+03 & 29.3 & 12.8 & 2.09e+09 & 4.58 & 21.0 & 8.3 & 1 \\
NC & 15 & 2.2e+02 & 29.3 & 12.8 & 2.86e+08 & 4.59 & 21.0 & 8.2 & 1 \\
NC & 15 & 3.1e+01 & 29.3 & 12.8 & 4.01e+07 & 4.63 & 21.1 & 8.2 & 1 \\
NC & 15 & 4.3e+00 & 29.3 & 12.8 & 5.59e+06 & 4.62 & 21.0 & 8.3 & 1 \\
NC & 15 & 6.0e-01 & 29.3 & 12.8 & 7.80e+05 & 4.57 & 21.1 & 8.2 & 1 \\
NC & 15 & 8.3e-02 & 29.2 & 12.7 & 1.09e+05 & 4.50 & 21.0 & 8.2 & 1 \\
NC & 15 & 1.2e-02 & 29.4 & 6.8 & 1.53e+04 & 4.51 & 21.0 & 8.4 & 1 \\
NC & 15 & 1.6e-03 & 31.0 & 6.4 & 1.02e+03 & 4.56 & 21.0 & 10.1 & 1 \\
T1 & 16 & 6.0e-01 & 37.5 & 7.7 & 1.14e+04 & 3.10 & 29.8 & 7.8 & 1 \\
T2 & 17 & 6.0e-01 & 37.5 & 7.9 & 1.18e+04 & 3.05 & 29.7 & 7.8 & 1 \\
T3 & 18 & 6.0e-01 & 44.0 & 7.2 & 1.32e+02 & 5.46 & 35.9 & 8.1 & 1 \\
T4 & 19 & 6.0e-01 & 43.9 & 11.8 & 1.19e+02 & 3.96 & 37.0 & 6.8 & 1 \\
\end{longtable}

\clearpage

\section*{References}\label{sec:references}
\addcontentsline{toc}{section}{References}

\protect\phantomsection\label{refs}
\begin{CSLReferences}{1}{0}
\bibitem[\citeproctext]{ref-kubo2002}
Kubo, T., Ohtani, E., Kato, T., Urakawa, S., Suzuki, A., Kanbe, Y., et
al. (2002). Mechanisms and kinetics of the post-spinel transformation in
{Mg\(_2\)SiO\(_4\)}. \emph{Physics of the Earth and Planetary
Interiors}, \emph{129}(1-2), 153--171.

\end{CSLReferences}


\end{document}
